\honest{Perpetual Motion Beliefs}{Eliezer Yudkowsky}{2008}

Yesterday I concluded that to form accurate beliefs about something ``you really do have to observe it,'' and that no clever argument can give you ``true knowledge of the unseen'' unless some step in it breaks the laws of physics. \nb{The notes on ``The Second Law of Thermodynamics, and Engines of Cognition'' discuss that post's claim that observing costs thermodynamic work.} The moral today: probability may be a state of belief, but its laws are ``harder than steel.''

People learn at school that certain beliefs are orders and probable ones mere suggestions. I state this without evidence. So the lottery buyer says: ``But you can't prove I won't win, right?''

Smash an egg. That it will not reassemble is only probable. But you must expect it to stay smashed: ``Probabilities are not certainties, but the laws of probability are theorems.''

I once knew a man whose spreadsheet showed that his wheels and gears produced reactionless thrust. Classical mechanics rules that out, so his spreadsheet had to contain a mistake. People who ``try to believe without evidence just this once'' build arguments like his machine, and ``there's always some step where a tiny probability turns into a large one.'' \nb{The post examines no such argument. Its one real case is the spreadsheet, a physics calculation.}

Why is the starting probability tiny? Because ``there is far more thin air than ground in the realms of Possibility.'' \nb{That is a claim about starting probabilities, and the laws of probability do not fix those; a subjective Bayesian may hold any coherent prior. In the egg, lottery and boiling-water cases, both sides agree the probability is tiny. For a belief held without evidence, that is the point in question.}

Stick your hand in boiling water, and you will be burned. Your ignorance of the water's exact state is what makes its energy heat. \nb{This is E.~T. Jaynes's view of heat, taken over from the previous post, and it is disputed.} Guessing the exact state is even less likely than being spared by chance. Each gear you add to the argument makes the machine less efficient, and each detail lowers the probability of the whole.

People ``really do construct these huge edifices of argument,'' I end, and should see them as that man's gears.
